\documentclass[11pt]{article}
\usepackage[margin=0.8in, top=0.5in]{geometry}
\usepackage{hyperref}
\usepackage{enumitem}

% Colors and Hyperlinks
\hypersetup{
    colorlinks=true,
    linkcolor=blue,
    urlcolor=blue,
    pdftitle={CV - Yichao Jin},
}

\begin{document}

% HEADER
\begin{center}
    {\huge \textbf{Yichao Jin}} \\
    \vspace{0.5em}
    School of Economic, Political and Policy Sciences \\
    The University of Texas at Dallas, Richardson, TX 75080 \\
    \vspace{0.5em}
    \href{https://yichao2022.github.io}{\textbf{https://yichao2022.github.io}} \quad | \quad \href{mailto:Yichao.Jin@UTDallas.edu}{Yichao.Jin@UTDallas.edu} \quad | \quad +1 (682) 266-4185 \\
    \vspace{0.5em}
    \textit{Specialization: Public Policy | Health Economics | Social Data Analytics | Behavioral Economics}
\end{center}

\vspace{1em}
\noindent\hrulefill

\section*{Education}
\textbf{University of Texas at Dallas} \hfill Expected Spring 2026 \\
\textbf{Ph.D., Public Policy and Political Economy} \\
- \textit{Distinction}: Recipient of the Dean of Graduate Education Dissertation Research Award (2025). \\
- \textit{Dissertation}: Dynamics of Intertemporal Choice in Vaccination Timing: A Discrete Choice Experiment Approach \\
- \textit{Chair}: Dohyeong Kim \quad | \quad \textit{Concentration}: Public Policy Analysis \& International Business

\vspace{0.5em}
\noindent \textbf{University of Texas at Dallas} \hfill Expected Spring 2026 \\
\textbf{M.S., Social Data Analytics and Research}

\vspace{0.5em}
\noindent \textbf{University of Queensland} \hfill 2019 \\
\textbf{Master of Development Economics}

\vspace{0.5em}
\noindent \textbf{University of California, Riverside} \hfill 2017 \\
\textbf{B.A., Economics}

\section*{Research Fields}
\textbf{Primary}: Public Policy, Health Economics, Applied Econometrics, Behavioral Economics \\
\textbf{Secondary}: Social Data Analytics, Policy Evaluation, Technology Governance (AI in Health)

\section*{Working Papers}
\begin{enumerate}
    \item \textbf{``How Humans Value Delayed Protection: Biological Vulnerability and Intertemporal Health Decisions''} (with Dohyeong Kim, Zhen Tian) \\
    \textit{Job Market Paper} --- \href{https://yichao2022.github.io/jmp.pdf}{\textbf{Link to Draft}}
    \begin{itemize}[noitemsep, topsep=0pt]
        \item Developed a structural Discrete Choice Experiment (DCE) framework to recover formal time-preference parameters (hyperbolic $\kappa = 0.225$, exponential $\delta = 0.160$) for delayed protection.
        \item Integrated behavioral parameters into dynamic SEIR models, identifying policy-critical gaps in vaccination uptake and resource allocation.
    \end{itemize}

    \item \textbf{``How CEO Succession Influences Firm AI Innovation: The Role of Academic Experience''} (with Zhimin Tian, Yu Xiang) \\
    \textit{Under Review at Journal of Business Research}

    \item \textbf{``The Ghost in the Machine: Algorithmic Rationality vs. Social Legitimacy in AI-Driven Public Health Governance''} \\
    \textit{Working Paper} --- Analyzes the theoretical tension between automated efficiency and institutional trust using the AGIL framework and Rational Choice models.

    \item \textbf{``The Price of Waiting: Evidence on Cash–Time Trade-Offs in Vaccination Time Discount''} \\
    \textit{In Preparation}
\end{enumerate}

\section*{Research Experience \& Affiliations}
\textbf{Doctoral Researcher, Spatial Health AI Research Partnership (SHARP), UT Dallas} \\
- Leading interdisciplinary research on the intersection of artificial intelligence, geospatial analytics, and health policy. \\
- Developing computational simulations (R, Python) to evaluate pandemic preparedness and public health administration.

\vspace{0.5em}
\noindent \textbf{Graduate Researcher, Behavioral Health Economics Laboratory, UT Dallas} \\
- Conducted large-scale behavioral experiments (N=1,027) analyzing risk-weighted health decisions and institutional trust.

\section*{Teaching Experience}
\textbf{Teaching Assistant, School of Economic, Political and Policy Sciences, UT Dallas}
\begin{itemize}[noitemsep, topsep=0pt]
    \item \textbf{Quantitative Methods for Policy Analysis (Graduate)}: Specialized in R programming, regression analysis, and causal inference for administrative data.
    \item \textbf{Health Economics \& Public Policy (Undergraduate)}: Focused on health market failures and the economics of preventive care.
    \item \textbf{Public Policy Analysis (Undergraduate)}: Mentored students in evidence-based policy memos and program evaluation.
\end{itemize}

\section*{Honors \& Awards}
\textbf{2025} \quad Dean of Graduate Education Dissertation Research Award, UT Dallas \\
\textbf{2025} \quad Betty \& Gifford Johnson Graduate Travel Award, UT Dallas \\
\textbf{2025} \quad Stanford Professional Certificate in Artificial Intelligence (AI Foundations, Machine Learning)

\section*{Technical Skills}
\textbf{Econometrics}: Mixed Logit (MXL), Latent Class Models (LCM), SEIR Modeling, Causal Inference. \\
\textbf{Software}: R, Stata (Advanced), Python, MATLAB, LaTeX, SQL, GIS. \\
\textbf{Languages}: English (Fluent), Chinese (Native).

\section*{Professional References}
\textbf{Dohyeong Kim (Chair)} | \textit{dohyeong.kim@utdallas.edu} \\
\textbf{Richard Scotch} | \textit{richard.scotch@utdallas.edu} \\
\textbf{Soyoung Kwon} | \textit{soyoung.kwon@utdallas.edu}

\vspace{2em}
\begin{center}
    \textit{Last Updated: March 7, 2026}
\end{center}

\end{document}
